\documentclass[conference]{IEEEtran}
\usepackage{booktabs}
\usepackage{multirow}
\usepackage{array}
\usepackage{adjustbox}
\usepackage{listings}
\lstset{
  language=Python,
  basicstyle=\ttfamily\scriptsize,
  keywordstyle=\bfseries,
  breaklines=true,
  frame=single,
  xleftmargin=1em,
  xrightmargin=1em,
  aboveskip=0.5em,
  belowskip=0.5em,
}

\title{Comparison of Results on the Gowalla Dataset}
\author{Compiled Report}
\date{}

\begin{document}
\maketitle

\section{POI Category Classification [Texas]}



\begin{table}[htbp]
\centering
\caption{F1-Score (\%) for POI Category Classification on Gowalla.}
\label{tab:classification}
\adjustbox{max width=\columnwidth}{
\begin{tabular}{lccccc}
\toprule
\textbf{Category} & \textbf{PGC-NN} & \textbf{MTL} & \textbf{HAVANA} & \textbf{ST-MTLNet} & \textbf{MTL\_new} \\
\midrule
Community     & 56.5$\pm$0.7  & \textit{57.37$\pm$0.69} & 29.46$\pm$1.7 & \underline{68.26$\pm$0.49} & \textbf{72.98$\pm$0.52} \\
Entertainment & 60.0$\pm$1.8  & 46.69$\pm$0.46 & \underline{67.02$\pm$0.5} & \textit{65.59$\pm$0.57} & \textbf{68.92$\pm$0.60} \\
Food          & 59.2$\pm$1.0  & 54.31$\pm$0.42 & \textbf{72.56$\pm$0.4} & \textit{66.50$\pm$0.24} & \underline{70.37$\pm$0.25} \\
Nightlife     & 54.6$\pm$3.1  & 34.44$\pm$1.21 & \textbf{77.95$\pm$0.7} & \textit{62.99$\pm$0.71} & \underline{66.56$\pm$0.75} \\
Outdoors      & 4.3$\pm$1.1   & 45.99$\pm$1.70 & \textit{54.08$\pm$0.8} & \underline{54.61$\pm$0.81} & \textbf{58.55$\pm$0.87} \\
Shopping      & 42.2$\pm$1.7  & \textbf{62.70$\pm$0.29} & 50.77$\pm$0.9 & \textit{57.10$\pm$0.18} & \underline{61.11$\pm$0.19} \\
Travel        & \textit{67.8$\pm$1.5}  & 39.37$\pm$1.29 & 66.49$\pm$1.3 & \underline{67.97$\pm$0.28} & \textbf{73.19$\pm$0.30} \\
\bottomrule
\end{tabular}
}
\end{table}

\section{Next POI Prediction [Texas]}

\begin{table}[htbp]
\centering
\caption{F1-Score (\%) for Next POI Prediction on Gowalla.}
\label{tab:nextpoi}
\adjustbox{max width=\columnwidth}{
\begin{tabular}{lcccc}
\toprule
\textbf{Category} & \textbf{POI-RGNN} & \textbf{MTL} & \textbf{ST-MTLNet} & \textbf{MTL\_new} \\
\midrule
Community     & \textit{35.7$\pm$0.5}  & 34.22$\pm$0.43 & \underline{40.21$\pm$1.25} & \textbf{42.33$\pm$1.32} \\
Entertainment & 22.7$\pm$0.6  & \textit{25.56$\pm$1.56} & \underline{26.63$\pm$1.46} & \textbf{28.30$\pm$1.55} \\
Food          & \underline{50.3$\pm$0.5}  & 29.56$\pm$4.10 & \textit{48.42$\pm$0.66} & \textbf{50.88$\pm$0.69} \\
Nightlife     & \underline{27.9$\pm$0.8}  & 25.46$\pm$1.13 & \textit{27.25$\pm$1.64} & \textbf{28.79$\pm$1.73} \\
Outdoors      & \underline{26.0$\pm$1.3}  & 23.02$\pm$1.03 & \textit{24.71$\pm$1.07} & \textbf{26.32$\pm$1.14} \\
Shopping      & \textbf{45.5$\pm$0.4}  & 38.79$\pm$4.12 & \textit{41.58$\pm$1.53} & \underline{43.69$\pm$1.61} \\
Travel        & 18.5$\pm$0.6  & \textit{29.71$\pm$1.32} & \underline{32.59$\pm$1.01} & \textbf{34.41$\pm$1.07} \\
\bottomrule
\end{tabular}
}
\end{table}

\section{News on MTLNet}

The improvements observed in MTL\_new stem from two key architectural changes over the original MTLnet.
\\
\\

\textbf{Category Head:} We replaced the simple MLP classifier (\texttt{CategoryHeadSingle}) with \texttt{CategoryHeadResidual}, which introduces skip connections between layers. Residual connections allow the network to learn refinements on top of the identity mapping, mitigating vanishing gradients in deeper heads and enabling better feature reuse. This is particularly beneficial when the input embedding already carries a rich representation from the GNN encoder, the residual path preserves that information while the learned path refines it for classification.
\newpage
\begin{lstlisting}[caption={ResidualBlock: skip connection preserves the original signal.}]
class ResidualBlock(nn.Module):
    def __init__(self, dim, hidden_dim, output_dim, dropout):
        super().__init__()
        self.net = nn.Sequential(
            nn.LayerNorm(dim),
            nn.Linear(dim, hidden_dim),   # expand
            nn.GELU(),
            nn.Dropout(dropout),
            nn.Linear(hidden_dim, output_dim),  # project
            nn.Dropout(dropout),
        )
        # match dimensions when they change
        self.shortcut = nn.Linear(dim, output_dim)

    def forward(self, x):
        return self.net(x) + self.shortcut(x)  # residual
\end{lstlisting}

\textbf{Embedding Fusion:} Instead of summing the spatial and temporal embeddings element-wise, we concatenate them into a single 128-dimensional vector. Summing two embeddings that originate from fundamentally different representation spaces (e.g., a graph-based spatial embedding and a Time2Vec temporal embedding) is semantically incorrect: it conflates dimensions that encode unrelated properties, destroying information and introducing noise. By concatenating, we preserve the full information from both sources and delegate the fusion to the model itself. In particular, the shared layers employ a FiLM (Feature-wise Linear Modulation) layer that learns per-dimension scaling ($\gamma$) and shifting ($\beta$) conditioned on a task embedding. Combined with the subsequent residual blocks, this allows the network's neurons to selectively attend to the dimensions of each embedding that are most informative for each task, effectively learning \emph{which} parts of each embedding are useful rather than assuming a na\"{i}ve one-to-one correspondence through summation.

\end{document}
